\documentclass{article}
\usepackage{amsmath,amssymb,amsthm,mathtools}
\usepackage[vocab=unitsquare]{leanify}
\newtheorem{proposition}{Proposition}
\newcommand{\CC}{\mathcal{C}}
\title{The unit interval as a type: points, integrals and measures of the field \emph{unitsquare}}
% A test of the type `[0,1]` that the field `unitsquare` declares: variables in `[0,1]`,
% functions on the square, numbers used as points (`C(\frac12,\frac12)`, `C(t,1-t)`),
% the integral `\int_0^1` and iterated integrals, Lebesgue measure, Borel measurable maps
% and the law of a pair. The expected output of
% `leanify check --vocab unitsquare --claims --defined` is check.txt: the readings, and
% what each claim and definition needs in order to be meaningful.
\begin{document}
\maketitle

An article that names the field \emph{unitsquare} reads $[0,1]$ as a type of its own: a variable in $[0,1]$ is a point of it, and a function on $[0,1]^2$ has no values outside the square.
A number used as a point, such as $1-t$, has to lie in $[0,1]$, and the integral $\int_0^1$ is the integral over the type.
Other intervals are sets of real numbers.

\section{Notions}

Let $\CC$ denote the class of bivariate copulas.
Lebesgue measure on $[0,1]$ is denoted by $\mathrm{Leb}$.
For $u,v\in[0,1]$, set $M(u,v)\coloneqq\min\{u,v\}$, $W(u,v)\coloneqq\max\{u+v-1,0\}$ and $\Pi(u,v)\coloneqq uv$.
For $s\in[0,1]$ and $u,v\in[0,1]$, set $S_s(u,v)\coloneqq\min\{u,v,1-s\}+\max\{u+v+s-2,0\}$.
We call a copula $C$ \emph{positively quadrant dependent} (\emph{PQD}) if $C(u,v)\ge uv$ for all $u,v\in[0,1]$.
A copula $C$ is called \emph{left-tail decreasing} (\emph{LTD}) if, for every $v\in[0,1]$, the map
\begin{equation}\label{eq:ltd}
u\longmapsto\frac{C(u,v)}{u},\qquad u\in(0,1],
\end{equation}
is non-increasing.
Define $\rho:\CC\to\mathbb{R}$ by $\rho(C)\coloneqq 12\int_0^1\int_0^1 C(u,v)\,du\,dv-3$.
Define $\beta:\CC\to\mathbb{R}$ by $\beta(C)\coloneqq 4C(\frac12,\frac12)-1$.
Define $\phi:\CC\to\mathbb{R}$ by $\phi(C)\coloneqq 6\int_0^1 C(t,t)\,dt-2$.
Define $\delta:\CC\to\mathbb{R}$ by $\delta(C)\coloneqq \int_0^1 C(t,1-t)\,dt$.
For a Borel measurable $T:[0,1]\to[0,1]$, let $\gamma_T$ denote the law of $(X,T(X))$ for $X$ uniformly distributed on $[0,1]$.
We say that a copula $F$ is \emph{supported on the graph} of a Borel measurable $f:[0,1]\to[0,1]$ if $F(u,v)=\mathrm{Leb}(\{x\in[0,1] : x\le u,\ f(x)\le v\})$ for all $u,v\in[0,1]$.

\section{Results}

\begin{proposition}\label{p1}
Every LTD copula is PQD.
\end{proposition}

\begin{proposition}\label{p2}
For every copula $C$, $-1\le\rho(C)\le 1$ and $\beta(C)\le 1$.
\end{proposition}

\begin{proposition}\label{p3}
Let $C$ be a copula. Then $W(u,v)\le C(u,v)\le M(u,v)$ for all $u,v\in[0,1]$.
\end{proposition}

\begin{proposition}\label{p4}
For every copula $C$ and every $u\in[0,1]$, $C(u,1-u)\le\frac12$ and $C(\frac u2,u)\le u$.
\end{proposition}

\begin{proposition}\label{p5}
The class $\CC$ is convex and compact with respect to uniform convergence. The function $\rho$ is continuous with respect to uniform convergence.
\end{proposition}

\begin{proposition}\label{p6}
For every $s\in[0,1]$, $S_s$ is a copula and $\phi(S_s)\le 1$. Moreover, $\rho(M)=1$, $\rho(W)=-1$ and $\rho(\Pi)=0$.
\end{proposition}

\begin{proposition}\label{p7}
$t\mapsto2t-1$ increases bijectively from $[0,1]$ onto $[-1,1]$.
\end{proposition}

\begin{proposition}\label{p8}
For every $u\in[0,1]$, $\int_0^1 M(u,t)\,dt=u-\frac{u^2}{2}$ and $\sup_{t\in[0,1]}\Pi(u,t)=u$.
\end{proposition}

\begin{proposition}\label{p9}
For every $x\in[0,1]$ and every $y\in[0,2]$, $x+y\le 3$ and $\frac{y}{2}\in[0,1]$.
\end{proposition}

\begin{proposition}\label{p10}
For every copula $C$ and every $u\in[0,1]$, $\int_0^u C(t,t)\,dt\le u$ and $\int_0^1\int_0^1 C(s,t)\,\mathrm{d}s\,\mathrm{d}t\le 1$.
\end{proposition}

A function on the square has partial derivatives: in each variable, the derivative within $[0,1]$, which is one-sided at the end points.

\begin{proposition}\label{p11}
For every $u,v\in[0,1]$, $\partial_1\Pi(u,v)=v$ and $\partial_2\Pi(u,v)=u$.
Moreover, $\int_0^1\int_0^1\bigl(\partial_1\Pi(u,v)-v\bigr)^2\,du\,dv=0$.
\end{proposition}

Functions on the square are added, and multiplied by numbers, value by value. A letter
may name a function and, with a subscript, a family: $M$ and $M_r$ are different symbols.

For $r\in[0,1]$, set $M_r\coloneqq(1-r)\Pi+rM$ and $W_r\coloneqq(1-r)\Pi+rW$.
For a copula $C$, the diagonal mass of $C$ is $q(C)\coloneqq\int_0^1 C(t,t)\,dt$.

\begin{proposition}\label{p12}
For every $r\in[0,1]$ and $u,v\in[0,1]$, we have $W_r(u,v)\le M_r(u,v)\le M(u,v)$.
Moreover, $M_1=M$, $M_0=\Pi$ and $\frac{M+W}2\le M$.
\end{proposition}

A sentence may run through a list: its items are its statements, one after the other.
The value of a function may be named by the function's letter.

\begin{proposition}\label{p13}
Every copula $C$ with diagonal mass $q=q(C)$ satisfies
\begin{enumerate}[(i)]
\item $0\le q$, $q\le\frac12$, and $q(1-q)\le\frac14$,
\item $q^2\le q$.
\end{enumerate}
For every copula $C$ and all $u,v\in[0,1]$, we have $C(u,v)\le M(u,v)$.
\end{proposition}

The integral of a function on the square against a copula is its integral for the measure
of the copula. An extremum may be over the pairs of a square. A symbol may carry a star
or, while its letter alone names nothing, a numeral as superscript, and a family may have
two subscripts.

For a copula $C$, set $\tau(C)\coloneqq4\int_{[0,1]^2}C(u,v)\,\mathrm{d}C(u,v)-1$, $\kappa(C)\coloneqq4\max_{(u,v)\in[0,1]^2}|C(u,v)-uv|$ and $\Phi^2(C)\coloneqq90\int_0^1\int_0^1\bigl(C(u,v)-uv\bigr)^2\,du\,dv$.
For a copula $C$, set $\tau^*(C)\coloneqq\frac25\tau(C)+\frac4{15}\Phi^2(C)$.
For $r\in[0,1]$ and $g\in[-r,r]$, set $C_{r,g}\coloneqq(1-r)\Pi+\frac{r+g}2M+\frac{r-g}2W$.

\begin{proposition}\label{p14}
For every copula $C$, we have $|\tau(C)|\le1$, $\kappa(C)\le1$ and $\tau^*(C)\le\Phi^2(C)+1$.
For every $r\in[0,1]$, $C_{r,r}=M_r$ and $C_{r,-r}=W_r$.
Moreover, $\int_{[0,1]^2}uv\,\mathrm{d}M(u,v)=\frac13$ and $\sup_{u,v\in[0,1]}\Pi(u,v)=1$.
\end{proposition}

A class may be described by its members, and a supremum may be over the members of such
a class.

Let $\mathcal F$ denote the class of Borel measurable functions $f:[0,1]\to\mathbb{R}$ with $\int_0^1 f(x)\,dx=0$ and $\int_0^1 f(x)^2\,dx=1$.
For a copula $C$, set $R(C)\coloneqq\sup\{\int_{[0,1]^2} f(u)g(v)\,\mathrm{d}C(u,v) : f,g\in\mathcal F\}$.

\begin{proposition}\label{p15}
Every copula $C$ with $R=R(C)$ satisfies
\begin{enumerate}[(i)]
\item $|\beta(C)|\le R$ and $\kappa(C)\le R$,
\item $|\tau(C)|\le\frac13R(2+R)$,
\item $\Phi^2(C)\le R^2$.
\end{enumerate}
For every $r\in[0,1]$, we have $R(M_r)=r$ and $\tau(M_r)=\frac13r(2+r)$.
\end{proposition}

The inverse of a function with a subscript in its name is written like any inverse. The
integrand of an integral against a copula has to be defined at almost every point for
the measure of the copula, and a derivative at the point where it is taken.

For $x\in\mathbb{R}$, set $\Phi_{\mathrm N}(x)\coloneqq\frac1{\sqrt{2\pi}}\int_{-\infty}^xe^{-t^2/2}\,dt$.
For a copula $C$, set $\rho_{\mathrm N}(C)\coloneqq\int_{[0,1]^2}\Phi_{\mathrm N}^{-1}(u)\,\Phi_{\mathrm N}^{-1}(v)\,\mathrm{d}C(u,v)$ and $\xi(C)\coloneqq6\int_0^1\int_0^1\bigl(\partial_1C(u,v)-v\bigr)^2\,du\,dv$.

\begin{proposition}\label{p16}
Every copula $C$ with $R=R(C)$ satisfies $|\rho_{\mathrm N}(C)|\le R$ and $\xi(C)\le R^2$.
For every $r\in[0,1]$, we have $R(M_r)=R(W_r)=r$ and
\begin{gather*}
\rho_{\mathrm N}(M_r)=r,\qquad\xi(M_r)=\Phi^2(M_r)=r^2,\\
\rho_{\mathrm N}(W_r)=-r,\qquad\tau(W_r)=-\tfrac13r(2+r).
\end{gather*}
\end{proposition}

A function variable may have a class of the article as its domain, and a family may be
indexed by it. Braces around a pair are the set of that pair. What follows `with' after
a noun is said of the same object.

For a function $A:\CC\to\mathbb{R}$, let $\mathcal R_A\coloneqq\{(A(C),R(C)) : C\in\CC\}$.
For $t\in[0,1]$ and $u,v\in[0,1]$, let $Q_t(u,v)\coloneqq(1-t)uv+t\min\{u,v\}$.

\begin{proposition}\label{p17}
We have
\begin{align*}
\mathcal R_\beta=\mathcal R_\rho&=\bigl\{(m,r):\ |m|\le r\le1\bigr\},\\
\mathcal R_\kappa&=\{(0,0)\}\cup\bigl\{(k,r):\ 0<k\le r\le1\bigr\},\\
\mathcal R_{\tau^*}&=\{(0,0)\}\cup\bigl\{(s,r):\ 0<r\le1,\ 0<s\le\tfrac1{15}r^2(6+4r)\bigr\}.
\end{align*}
For every $t\in[0,1]$, $Q_t$ is a copula with $\beta(Q_t)=t$ and $R(Q_t)=t$.
\end{proposition}

A notation may have a subscript among its places: the expression it stands for is
meaningful where it is used, not everywhere.

We write $\ell_A(x)$ for $\min\{R(C):\ C\in\CC,\ A(C)=x\}$.

\begin{proposition}\label{p18}
For every $x\in[-1,1]$, we have $\ell_\beta(x)=\ell_\rho(x)=|x|$.
For every $x\in[0,1]$ and $r\in[0,1]$, we have $\ell_{\Phi^2}(x)=r$ if and only if $r^2=x$.
\end{proposition}
\end{document}
